\section*{Appendix A: Related Work}
\label{app:rel-work}

Multi-armed bandits \citep{slivkins2019introduction,lattimore2020bandit} have been studied extensively, which can be traced back to the foundational contributions of \citep{thompson1933likelihood,lai1985asymptotically}. This framework has proven remarkably versatile across a wide range of applications, including clinical trials \citep{gittins1}, recommendation systems \citep{Li_2010}, financial optimization \citep{brochu}, and modern incentivized learning methods \citep{fraz,wang,chakraborty2024incentivized,chakraborty2025incentivized}.

The literature on multi-agent bandits is vast; we highlight the main works introducing collisions, decentralization, and cooperative settings. In the discrete-arm case, \citep{rosenski2016multi} introduced the Musical Chairs protocol for decentralized coordination, later improved by \citep{wang2020optimal} with optimal regret guarantees. Variants with delayed feedback, fairness, or corruption have also been explored \citep{boursier2020survey,ghaffari2024robust}.

Our setting departs by considering a continuum of actions. Single-player Lipschitz bandits have been studied via discretization, zooming, and hierarchical methods \citep{magureanu2014lipschitz,bubeck2011xarmed,Kleinberg}. These works provide minimax rates but do not address multi-agent collisions.

Multi-agent bandits in continuous domains remain underexplored. Existing works often assume stronger structure (e.g., convexity, linearity) \citep{bambos}, or allow explicit communication \citep{landgren2020distributed}.

In contrast, to our knowledge, our work is the first to provide a fully decentralized, communication-free solution for multi-player Lipschitz bandits with hard collisions, achieving near-optimal regret with only a constant coordination overhead.
